% --- Archon physics formalization source begin ---
% archon:physics
% archon:covers PhyXMiniProblems/problem_phyx_mini_0861.lean
% archon:source-report reports/phyx_mini/problem_phyx_mini_0861.source.json

\chapter{Physics problem phyx\_mini\_0861}
\label{ch:PhyXMiniProblems_problem_phyx_mini_0861}

\paragraph{Problem source.}
\#\# Physical scenario

Two point charges \textbackslash{}( q\_a \textbackslash{}) and \textbackslash{}( q\_b \textbackslash{}) are located on the \textbackslash{}( x \textbackslash{})-axis at \textbackslash{}( x = a \textbackslash{}) and \textbackslash{}( x = b \textbackslash{}). figure is a graph of \textbackslash{}( E\_x \textbackslash{}), the \textbackslash{}( x \textbackslash{})-component of the electric field.

\#\# Question

What is the ratio \textbackslash{}( \textbackslash{}left| q\_a / q\_b \textbackslash{}right| \textbackslash{})?

\#\# Answer choices

- A: A: 1.2

- B: B: 1.4

- C: C: 2

- D: D: 1

\#\# Figure caption (auxiliary; use the image as primary evidence)

The image contains a mathematical graph featuring two curves. 

1. **Horizontal Line**: There's a solid black horizontal line labeled "x". 

2. **Vertical Lines**: Two dashed vertical lines intersect the horizontal line. They are labeled "a" and "b" at the points where they intersect the horizontal line.

3. **Curves**: 
   - Above the horizontal line, there's a magenta curve resembling an inverted parabola that is symmetric around the midpoint between "a" and "b".
   - Below the horizontal line, there is a magenta curve in the shape of a parabola, opening downwards.

These elements together create a symmetrical graph where the curves above and below the line are mirror images concerning their positioning relative to "a" and "b".

\#\# Dataset metadata

Electrostatics; Physical Model Grounding Reasoning, Multi-Formula Reasoning

\paragraph{Recorded answer/context.}
D: D: 1

\paragraph{Figure/image path.}
/root/proposal\_for\_physic/science-mango/phyx\_mini\_run/phyx\_data/test\_image/861.png

\paragraph{Formalization target.}
create a compiling Lean file with sorry bodies at `PhyXMiniProblems/problem\_phyx\_mini\_0861.lean`. The Lean declarations must preserve the physical quantities, dimensions or dimensional roles, figure labels, governing-law hypotheses, and final relation expressed by this problem.
Use Mathlib/Physlib names found through LeanExplore where available. If a physics API is missing, introduce faithful local abstractions rather than scalar placeholder aliases.

\begin{theorem}[Physics formalization target]
\label{thm:physics:phyx_mini_0861:target}
\lean{PhyXMiniProblems.ProblemPhyXMini0861.problem_phyx_mini_0861}
\uses{def:physics:phyx-mini-0861:phyxminiproblems-problemphyxmini0861-twopointchargeaxissetup, def:physics:phyx-mini-0861:phyxminiproblems-problemphyxmini0861-fieldcomponentrepresentsphyslibelectricfield, def:physics:phyx-mini-0861:phyxminiproblems-problemphyxmini0861-matchessuppliedelectricfieldgraph, def:physics:phyx-mini-0861:phyxminiproblems-problemphyxmini0861-hasphysicaltwopointchargeparameters, def:physics:phyx-mini-0861:phyxminiproblems-problemphyxmini0861-satisfiestwopointchargecoulomblaw, def:physics:phyx-mini-0861:phyxminiproblems-problemphyxmini0861-absolutechargeratio, def:physics:phyx-mini-0861:phyxminiproblems-problemphyxmini0861-recordeddatasetanswer, def:physics:phyx-mini-0861:phyxminiproblems-problemphyxmini0861-isuniquematchingdisplayedratio}
The assigned autoformalize agent should translate this physics problem into Lean declarations in the covered file, with theorem and lemma proof bodies written as `by sorry`.
\end{theorem}
\begin{proof}
This is an autoformalization task, not a proof task. Produce statements that can later be proved by the physics prover without weakening the physical model.
\end{proof}

% --- Archon declaration topology begin ---
\section{Declaration topology}
\begin{definition}[electric Field Strength Dimension]
  \label{def:physics:phyx-mini-0861:phyxminiproblems-problemphyxmini0861-electricfieldstrengthdimension}
  \lean{PhyXMiniProblems.ProblemPhyXMini0861.electricFieldStrengthDimension}
  The dimension M L T⁻² C⁻¹ of electric-field strength
\end{definition}
\begin{proof}
  This is the declared model definition of electric Field Strength Dimension.
\end{proof}
\begin{definition}[Axis Coordinate Quantity]
  \label{def:physics:phyx-mini-0861:phyxminiproblems-problemphyxmini0861-axiscoordinatequantity}
  \lean{PhyXMiniProblems.ProblemPhyXMini0861.AxisCoordinateQuantity}
  A signed, unit-independent coordinate on the physical x-axis
\end{definition}
\begin{proof}
  This is the declared model definition of Axis Coordinate Quantity.
\end{proof}
\begin{definition}[Signed Charge Quantity]
  \label{def:physics:phyx-mini-0861:phyxminiproblems-problemphyxmini0861-signedchargequantity}
  \lean{PhyXMiniProblems.ProblemPhyXMini0861.SignedChargeQuantity}
  A signed, unit-independent electric charge
\end{definition}
\begin{proof}
  This is the declared model definition of Signed Charge Quantity.
\end{proof}
\begin{definition}[Signed Electric Field Component Quantity]
  \label{def:physics:phyx-mini-0861:phyxminiproblems-problemphyxmini0861-signedelectricfieldcomponentquantity}
  \lean{PhyXMiniProblems.ProblemPhyXMini0861.SignedElectricFieldComponentQuantity}
  \uses{def:physics:phyx-mini-0861:phyxminiproblems-problemphyxmini0861-electricfieldstrengthdimension}
  A signed, unit-independent x-component of electric field
\end{definition}
\begin{proof}
  This is the declared model definition of Signed Electric Field Component Quantity.
\end{proof}
\begin{definition}[axis Coordinate In Meters]
  \label{def:physics:phyx-mini-0861:phyxminiproblems-problemphyxmini0861-axiscoordinateinmeters}
  \lean{PhyXMiniProblems.ProblemPhyXMini0861.axisCoordinateInMeters}
  \uses{def:physics:phyx-mini-0861:phyxminiproblems-problemphyxmini0861-axiscoordinatequantity}
  Coherent-SI coordinate readout, in metres
\end{definition}
\begin{proof}
  This is the declared model definition of axis Coordinate In Meters.
\end{proof}
\begin{definition}[charge In Coulombs]
  \label{def:physics:phyx-mini-0861:phyxminiproblems-problemphyxmini0861-chargeincoulombs}
  \lean{PhyXMiniProblems.ProblemPhyXMini0861.chargeInCoulombs}
  \uses{def:physics:phyx-mini-0861:phyxminiproblems-problemphyxmini0861-signedchargequantity}
  Coherent-SI signed-charge readout, in coulombs
\end{definition}
\begin{proof}
  This is the declared model definition of charge In Coulombs.
\end{proof}
\begin{definition}[field Component In Newtons Per Coulomb]
  \label{def:physics:phyx-mini-0861:phyxminiproblems-problemphyxmini0861-fieldcomponentinnewtonspercoulomb}
  \lean{PhyXMiniProblems.ProblemPhyXMini0861.fieldComponentInNewtonsPerCoulomb}
  \uses{def:physics:phyx-mini-0861:phyxminiproblems-problemphyxmini0861-signedelectricfieldcomponentquantity}
  Coherent-SI signed field-component readout, in newtons per coulomb
\end{definition}
\begin{proof}
  This is the declared model definition of field Component In Newtons Per Coulomb.
\end{proof}
\begin{definition}[Charge Label]
  \label{def:physics:phyx-mini-0861:phyxminiproblems-problemphyxmini0861-chargelabel}
  \lean{PhyXMiniProblems.ProblemPhyXMini0861.ChargeLabel}
  The two source labels printed below the dashed lines in the graph
\end{definition}
\begin{proof}
  This is the declared model definition of Charge Label.
\end{proof}
\begin{definition}[Axis Point Charge]
  \label{def:physics:phyx-mini-0861:phyxminiproblems-problemphyxmini0861-axispointcharge}
  \lean{PhyXMiniProblems.ProblemPhyXMini0861.AxisPointCharge}
  \uses{def:physics:phyx-mini-0861:phyxminiproblems-problemphyxmini0861-axiscoordinatequantity, def:physics:phyx-mini-0861:phyxminiproblems-problemphyxmini0861-signedchargequantity}
  A point charge together with its dimensionful x-axis coordinate
\end{definition}
\begin{proof}
  This is the declared model definition of Axis Point Charge.
\end{proof}
\begin{definition}[Field Graph Region]
  \label{def:physics:phyx-mini-0861:phyxminiproblems-problemphyxmini0861-fieldgraphregion}
  \lean{PhyXMiniProblems.ProblemPhyXMini0861.FieldGraphRegion}
  The three regions separated by the two vertical source lines
\end{definition}
\begin{proof}
  This is the declared model definition of Field Graph Region.
\end{proof}
\begin{definition}[Two Charge Electric Field Figure]
  \label{def:physics:phyx-mini-0861:phyxminiproblems-problemphyxmini0861-twochargeelectricfieldfigure}
  \lean{PhyXMiniProblems.ProblemPhyXMini0861.TwoChargeElectricFieldFigure}
  \uses{def:physics:phyx-mini-0861:phyxminiproblems-problemphyxmini0861-chargelabel, def:physics:phyx-mini-0861:phyxminiproblems-problemphyxmini0861-fieldgraphregion}
  Literal presentation features of image 861.png. Quantitative interpretation of the curve is stated separately below
\end{definition}
\begin{proof}
  This is the declared model definition of Two Charge Electric Field Figure.
\end{proof}
\begin{definition}[Two Point Charge Axis Setup]
  \label{def:physics:phyx-mini-0861:phyxminiproblems-problemphyxmini0861-twopointchargeaxissetup}
  \lean{PhyXMiniProblems.ProblemPhyXMini0861.TwoPointChargeAxisSetup}
  \uses{def:physics:phyx-mini-0861:phyxminiproblems-problemphyxmini0861-signedelectricfieldcomponentquantity, def:physics:phyx-mini-0861:phyxminiproblems-problemphyxmini0861-chargelabel, def:physics:phyx-mini-0861:phyxminiproblems-problemphyxmini0861-axispointcharge, def:physics:phyx-mini-0861:phyxminiproblems-problemphyxmini0861-twochargeelectricfieldfigure}
  Independent physical objects for the electrostatic configuration. The requested ratio is not stored here. The scalar component is a dimensionful observable indexed by an explicitly metre-valued coordinate, while electricField retains Physlib's full field role
\end{definition}
\begin{proof}
  This is the declared model definition of Two Point Charge Axis Setup.
\end{proof}
\begin{definition}[source Position In Meters]
  \label{def:physics:phyx-mini-0861:phyxminiproblems-problemphyxmini0861-sourcepositioninmeters}
  \lean{PhyXMiniProblems.ProblemPhyXMini0861.sourcePositionInMeters}
  \uses{def:physics:phyx-mini-0861:phyxminiproblems-problemphyxmini0861-axiscoordinateinmeters, def:physics:phyx-mini-0861:phyxminiproblems-problemphyxmini0861-chargelabel, def:physics:phyx-mini-0861:phyxminiproblems-problemphyxmini0861-twopointchargeaxissetup}
  Metre coordinate of a named point charge
\end{definition}
\begin{proof}
  This is the declared model definition of source Position In Meters.
\end{proof}
\begin{definition}[source Charge In Coulombs]
  \label{def:physics:phyx-mini-0861:phyxminiproblems-problemphyxmini0861-sourcechargeincoulombs}
  \lean{PhyXMiniProblems.ProblemPhyXMini0861.sourceChargeInCoulombs}
  \uses{def:physics:phyx-mini-0861:phyxminiproblems-problemphyxmini0861-chargeincoulombs, def:physics:phyx-mini-0861:phyxminiproblems-problemphyxmini0861-chargelabel, def:physics:phyx-mini-0861:phyxminiproblems-problemphyxmini0861-twopointchargeaxissetup}
  Coulomb readout of a named point charge
\end{definition}
\begin{proof}
  This is the declared model definition of source Charge In Coulombs.
\end{proof}
\begin{definition}[field XIn Newtons Per Coulomb]
  \label{def:physics:phyx-mini-0861:phyxminiproblems-problemphyxmini0861-fieldxinnewtonspercoulomb}
  \lean{PhyXMiniProblems.ProblemPhyXMini0861.fieldXInNewtonsPerCoulomb}
  \uses{def:physics:phyx-mini-0861:phyxminiproblems-problemphyxmini0861-fieldcomponentinnewtonspercoulomb, def:physics:phyx-mini-0861:phyxminiproblems-problemphyxmini0861-twopointchargeaxissetup}
  Signed graph ordinate E\_x, in newtons per coulomb
\end{definition}
\begin{proof}
  This is the declared model definition of field XIn Newtons Per Coulomb.
\end{proof}
\begin{definition}[reflect About Source Midpoint]
  \label{def:physics:phyx-mini-0861:phyxminiproblems-problemphyxmini0861-reflectaboutsourcemidpoint}
  \lean{PhyXMiniProblems.ProblemPhyXMini0861.reflectAboutSourceMidpoint}
  \uses{def:physics:phyx-mini-0861:phyxminiproblems-problemphyxmini0861-twopointchargeaxissetup, def:physics:phyx-mini-0861:phyxminiproblems-problemphyxmini0861-sourcepositioninmeters}
  Reflection of an x-coordinate about the midpoint of a and b
\end{definition}
\begin{proof}
  This is the declared model definition of reflect About Source Midpoint.
\end{proof}
\begin{definition}[Field Component Represents Physlib Electric Field]
  \label{def:physics:phyx-mini-0861:phyxminiproblems-problemphyxmini0861-fieldcomponentrepresentsphyslibelectricfield}
  \lean{PhyXMiniProblems.ProblemPhyXMini0861.FieldComponentRepresentsPhyslibElectricField}
  \uses{def:physics:phyx-mini-0861:phyxminiproblems-problemphyxmini0861-twopointchargeaxissetup, def:physics:phyx-mini-0861:phyxminiproblems-problemphyxmini0861-fieldxinnewtonspercoulomb}
  The scalar ordinate used by the graph is the first component of Physlib's one-dimensional electric field at the corresponding axis point. Both sides are coherent-SI readouts
\end{definition}
\begin{proof}
  This is the declared model definition of Field Component Represents Physlib Electric Field.
\end{proof}
\begin{definition}[Matches Supplied Electric Field Graph]
  \label{def:physics:phyx-mini-0861:phyxminiproblems-problemphyxmini0861-matchessuppliedelectricfieldgraph}
  \lean{PhyXMiniProblems.ProblemPhyXMini0861.MatchesSuppliedElectricFieldGraph}
  \uses{def:physics:phyx-mini-0861:phyxminiproblems-problemphyxmini0861-twopointchargeaxissetup, def:physics:phyx-mini-0861:phyxminiproblems-problemphyxmini0861-sourcepositioninmeters, def:physics:phyx-mini-0861:phyxminiproblems-problemphyxmini0861-fieldxinnewtonspercoulomb, def:physics:phyx-mini-0861:phyxminiproblems-problemphyxmini0861-reflectaboutsourcemidpoint}
  Literal and interpreted evidence from the primary image. In particular, curveMirrorSymmetry records a visible feature of the supplied graph; it does not directly assert any relation between the two charges
\end{definition}
\begin{proof}
  This is the declared model definition of Matches Supplied Electric Field Graph.
\end{proof}
\begin{definition}[Has Physical Two Point Charge Parameters]
  \label{def:physics:phyx-mini-0861:phyxminiproblems-problemphyxmini0861-hasphysicaltwopointchargeparameters}
  \lean{PhyXMiniProblems.ProblemPhyXMini0861.HasPhysicalTwoPointChargeParameters}
  \uses{def:physics:phyx-mini-0861:phyxminiproblems-problemphyxmini0861-twopointchargeaxissetup, def:physics:phyx-mini-0861:phyxminiproblems-problemphyxmini0861-sourcechargeincoulombs}
  Nondegeneracy and positivity conditions for two genuine point charges
\end{definition}
\begin{proof}
  This is the declared model definition of Has Physical Two Point Charge Parameters.
\end{proof}
\begin{definition}[Satisfies Two Point Charge Coulomb Law]
  \label{def:physics:phyx-mini-0861:phyxminiproblems-problemphyxmini0861-satisfiestwopointchargecoulomblaw}
  \lean{PhyXMiniProblems.ProblemPhyXMini0861.SatisfiesTwoPointChargeCoulombLaw}
  \uses{def:physics:phyx-mini-0861:phyxminiproblems-problemphyxmini0861-twopointchargeaxissetup, def:physics:phyx-mini-0861:phyxminiproblems-problemphyxmini0861-sourcepositioninmeters, def:physics:phyx-mini-0861:phyxminiproblems-problemphyxmini0861-sourcechargeincoulombs, def:physics:phyx-mini-0861:phyxminiproblems-problemphyxmini0861-fieldxinnewtonspercoulomb}
  The one-dimensional point-charge law and linear superposition. Away from both source positions, each term is k q r / |r|³; their sum is the measured resultant x-component. This all-regular-points law contains no requested charge ratio
\end{definition}
\begin{proof}
  This is the declared model definition of Satisfies Two Point Charge Coulomb Law.
\end{proof}
\begin{definition}[absolute Charge Ratio]
  \label{def:physics:phyx-mini-0861:phyxminiproblems-problemphyxmini0861-absolutechargeratio}
  \lean{PhyXMiniProblems.ProblemPhyXMini0861.absoluteChargeRatio}
  \uses{def:physics:phyx-mini-0861:phyxminiproblems-problemphyxmini0861-twopointchargeaxissetup, def:physics:phyx-mini-0861:phyxminiproblems-problemphyxmini0861-sourcechargeincoulombs}
  The dimensionless magnitude ratio asked for in the problem
\end{definition}
\begin{proof}
  This is the declared model definition of absolute Charge Ratio.
\end{proof}
\begin{definition}[Answer Choice]
  \label{def:physics:phyx-mini-0861:phyxminiproblems-problemphyxmini0861-answerchoice}
  \lean{PhyXMiniProblems.ProblemPhyXMini0861.AnswerChoice}
  The four answer labels supplied with the problem
\end{definition}
\begin{proof}
  This is the declared model definition of Answer Choice.
\end{proof}
\begin{definition}[displayed Charge Magnitude Ratio]
  \label{def:physics:phyx-mini-0861:phyxminiproblems-problemphyxmini0861-answerchoice-displayedchargemagnituderatio}
  \lean{PhyXMiniProblems.ProblemPhyXMini0861.AnswerChoice.displayedChargeMagnitudeRatio}
  \uses{def:physics:phyx-mini-0861:phyxminiproblems-problemphyxmini0861-answerchoice}
  Exact real value printed by each displayed answer choice
\end{definition}
\begin{proof}
  This is the declared model definition of displayed Charge Magnitude Ratio.
\end{proof}
\begin{definition}[recorded Dataset Answer]
  \label{def:physics:phyx-mini-0861:phyxminiproblems-problemphyxmini0861-recordeddatasetanswer}
  \lean{PhyXMiniProblems.ProblemPhyXMini0861.recordedDatasetAnswer}
  \uses{def:physics:phyx-mini-0861:phyxminiproblems-problemphyxmini0861-answerchoice}
  The answer label recorded in the supplied dataset metadata
\end{definition}
\begin{proof}
  This is the declared model definition of recorded Dataset Answer.
\end{proof}
\begin{definition}[Is Unique Matching Displayed Ratio]
  \label{def:physics:phyx-mini-0861:phyxminiproblems-problemphyxmini0861-isuniquematchingdisplayedratio}
  \lean{PhyXMiniProblems.ProblemPhyXMini0861.IsUniqueMatchingDisplayedRatio}
  \uses{def:physics:phyx-mini-0861:phyxminiproblems-problemphyxmini0861-answerchoice}
  A displayed choice is the unique one whose value equals a given ratio
\end{definition}
\begin{proof}
  This is the declared model definition of Is Unique Matching Displayed Ratio.
\end{proof}
\begin{lemma}[source Charge Readouts sum eq zero]
  \label{lem:physics:phyx-mini-0861:phyxminiproblems-problemphyxmini0861-sourcechargereadouts-sum-eq-zero}
  \lean{PhyXMiniProblems.ProblemPhyXMini0861.sourceChargeReadouts_sum_eq_zero}
  \uses{def:physics:phyx-mini-0861:phyxminiproblems-problemphyxmini0861-twopointchargeaxissetup, def:physics:phyx-mini-0861:phyxminiproblems-problemphyxmini0861-sourcechargeincoulombs, def:physics:phyx-mini-0861:phyxminiproblems-problemphyxmini0861-matchessuppliedelectricfieldgraph, def:physics:phyx-mini-0861:phyxminiproblems-problemphyxmini0861-hasphysicaltwopointchargeparameters, def:physics:phyx-mini-0861:phyxminiproblems-problemphyxmini0861-satisfiestwopointchargecoulomblaw}
  Midpoint-reflection symmetry of the Coulomb field forces the signed source charges to be opposites. This is a derived lemma, not a premise
\end{lemma}
\begin{proof}
  The conclusion follows from the governing relations and figure readouts named in the statement of source Charge Readouts sum eq zero.
\end{proof}
% --- Archon declaration topology end ---
% --- Archon physics formalization source end ---
